\chapter{Diferenciación en espacios normados}
\label{ch:diferenciacion}
\index{Diferencial de Fréchet}\index{Derivada de Fréchet}\index{Teorema de la función inversa}\index{Regla de la cadena}

\section{El diferencial de Fréchet}

La noción de derivada se extiende de manera natural del cálculo en una variable
al contexto de espacios normados. La idea central sigue siendo la aproximación
lineal local de una función.

\begin{definition}\label{def:frechet}
Sean $E$ y $F$ espacios normados, $U\subseteq E$ abierto y $f\colon U\to F$.
Se dice que $f$ es \emph{diferenciable} (en el sentido de Fréchet) en
$x_0\in U$ si existe una aplicación lineal continua $T\colon E\to F$ tal que
\begin{equation}\label{eq:frechet}
  \lim_{h\to 0}\frac{\|f(x_0+h)-f(x_0)-T(h)\|}{\|h\|}=0.
\end{equation}
En tal caso, $T$ es única y se denomina \emph{diferencial} de $f$ en $x_0$,
notada $Df(x_0)$ o $f'(x_0)$.
\end{definition}

\begin{remark}
Cuando $E=F=\mathbb{R}$, la definición anterior coincide con la derivada usual:
$T(h)=f'(x_0)\,h$.
\end{remark}

\begin{example}\label{ej:derivada-lineal}
Si $f\colon E\to F$ es lineal y continua, entonces $f$ es diferenciable en
todo punto y $Df(x)=f$ para todo $x\in E$.
\end{example}

\begin{example}\label{ej:derivada-bilineal}
Sea $b\colon E_1\times E_2\to F$ una aplicación bilineal continua. Entonces $b$
es diferenciable en todo punto y su diferencial está dada por
\[
  Db(x_1,x_2)(h_1,h_2)=b(x_1,h_2)+b(h_1,x_2).
\]
\end{example}

\section{Reglas de diferenciación}

\begin{proposition}[Regla de la cadena]\label{prop:cadena}
Sean $U\subseteq E$ abierto, $V\subseteq F$ abierto, $f\colon U\to V$
diferenciable en $x_0$ y $g\colon V\to G$ diferenciable en $f(x_0)$. Entonces
$g\circ f$ es diferenciable en $x_0$ y
\begin{equation}\label{eq:cadena}
  D(g\circ f)(x_0)=Dg(f(x_0))\circ Df(x_0).
\end{equation}
\end{proposition}

\begin{theorem}[del valor medio]\label{teo:valor-medio}
Sea $f\colon U\to F$ diferenciable en el segmento $[x,y]\subseteq U$. Entonces
\begin{equation}\label{eq:valor-medio}
  \|f(y)-f(x)\|\leq\sup_{z\in[x,y]}\|Df(z)\|\,\|y-x\|.
\end{equation}
\end{theorem}

\section{El teorema de la función inversa}

\begin{theorem}[de la función inversa]\label{teo:inversa}
Sea $f\colon U\subseteq E\to F$ de clase $C^1$ (es decir, continuamente
diferenciable) y sea $x_0\in U$ tal que $Df(x_0)$ es un isomorfismo lineal
sobreyectivo. Entonces existen entornos abiertos $V\subseteq U$ de $x_0$ y
$W\subseteq F$ de $f(x_0)$ tales que $f\colon V\to W$ es biyectiva y su
inversa $f^{-1}\colon W\to V$ es de clase $C^1$. Además,
\begin{equation}\label{eq:inversa}
  D(f^{-1})(f(x_0))=\bigl(Df(x_0)\bigr)^{-1}.
\end{equation}
\end{theorem}

Este teorema constituye uno de los pilares del cálculo diferencial en varias
variables y en dimensiones infinitas, pues garantiza que la linealización local
de una función controla su comportamiento local no lineal siempre que la
linealización sea invertible.

\section{Ejercicios}

\begin{exercise}
Calcule el diferencial de Fréchet de la función $f\colon\mathbb{R}^2\to\mathbb{R}$
definida por $f(x,y)=x^2+xy+y^2$ en un punto genérico $(x_0,y_0)$.
\end{exercise}

\begin{exercise}
Demuestre que si $f$ es diferenciable en $x_0$, entonces $f$ es continua en
$x_0$.
\end{exercise}

\begin{exercise}
Sea $f\colon M_n(\mathbb{R})\to M_n(\mathbb{R})$ dada por $f(A)=A^2$. Calcule
$Df(A)(H)$ para matrices $A,H\in M_n(\mathbb{R})$.
\end{exercise}

\begin{exercise}
Use el \cref{teo:inversa} para probar que existe un entorno de la matriz
identidad $I$ en $M_n(\mathbb{R})$ en el cual toda matriz posee raíz cuadrada.
\end{exercise}

\begin{problem}
Formule y demuestre un teorema de la función implícita en espacios de Banach
usando el \cref{teo:inversa} como herramienta principal.
\end{problem}
